\subsection{Jointness measures}\label{sec:jointness_measures}
The main implementations of Bayesian model averaging are concerned
with the selection of variables when model uncertainty is
present. Another relevant issue which arises in this framework is to
identify whether different sets of two variables $x_i$ and $x_j$ are
substitutes, complements or neither over the model space. For that
reason, \cite{LeySteel2007} and \cite{DoppelhoferWeeks2009} define
ex-post jointness measures of dependence between different sets of
explanatory variables. Following \cite{BivariateJointness-2018},
let's define \nameref{tab:GeneralContingencyTable}:

\begin{table}[htb]
	\centering
	\caption{General contingency table for two binary variables.}
	\begin{tabular}{l|ccc}\hline
						&\textit{B}	&	$\overline{B}$	&Total	\\\hline
		\textit{A}		&	a		&	b			& $a+b$		\\
		$\overline{A}$	&	c		&	d			& $c+d$		\\
		Total			&	$a+c$	&	$b+d$		& $a+b+c+d=N$\\\hline
	\end{tabular}
	\label{tab:GeneralContingencyTable}
\end{table}

\begin{enumerate}
	\item The jointness measure was proposed for the first
		time in \cite{DoppelhoferWeeks2005} as
		(\lstinline!"Doppelhofer-Weeks Measure (DW1)"!
		in our package):
		\begin{equation}
			J_{\mathit{DW1}}(A,B) = \ln\left(\frac{\Prob(AB\mid y)}
			{\Prob(A\mid y) \Prob(B\mid y)}\right)\sim
			J_{\mathit{DW1}} = \ln\left(\frac{aN}{(a+b)(a+c)}\right).
			\label{DW_jointness1}
		\end{equation}
	\item An alternative jointness measure was proposed by
		\cite{DoppelhoferWeeks2009} as
		(\lstinline!"Doppelhofer-Weeks Measure (DW2)"!
		in our package):
		\begin{equation}
			J_{\mathit{DW2}}(A,B) =
			\ln\left(\frac {\Prob(AB \mid y)\Prob(\overline{A}\overline{B} \mid y)}
			{\Prob(A\overline{B}\mid y)\Prob(\overline{A}B\mid y)}\right)\sim
			J_{\mathit{DW2}} = \ln\left(\frac{ad}{bc}\right).
			\label{DW_jointness2}
		\end{equation}
		%where $\Prob(\tilde{i} \cap \tilde{j}\mid y)$ denotes the sum of the
		%posterior probabilities of the regression models in which neither
		%$x_i$ and $x_j$ are included, $\Prob(i \cap \tilde{j}\mid y)$
		%corresponds to the sum of the posterior probabilities of all models in
		%which $x_i$ is included and $x_j$ is excluded. The last probability
		%$\Prob(\tilde{i} \cap j\mid y)$ is defined analogously.
		%$J_{\mathit{DW}}$ corresponds to the posterior odds of including $i$,
		%given that $j$ is included, divided by the posterior odds of including
		%$i$ given that $j$ is not included (see %\citealt{DoppelhoferWeeks2009}).,
	\item Two jointness statistics were proposed
		in \cite{LeySteel2007} as
		(\lstinline!"Ley-Steel Measure (LS1)"! in our package):
		\begin{equation}
			J_{\mathit{LS1}}(A,B) =
			\frac{\Prob(AB\mid y)}{\Prob(A\mid y)+\Prob(B\mid y)-\Prob(AB\mid y)}
			\sim J_{\mathit{LS1}} = \frac{a}{a+b+c},
			\label{LS_jointness1}
		\end{equation}
	\item and (\lstinline!"Ley-Steel Measure (LS2)"! in our package):
	\begin{equation}
		J_{\mathit{LS2}}(A,B) =
		\frac{\Prob(AB\mid y)}{\Prob(A\overline{B}\mid y)+\Prob(\overline{A}B\mid y)}
		\sim J_{\mathit{LS2}} = \frac{a}{b+c}.
		\label{LS_jointness2}
	\end{equation}
%	where $\Prob(i \cap j\mid y)$ represents the sum of the posterior
%	probabilities of those models that contain both variables $x_i$ and
%	$x_j$, $\Prob(i\mid y)$ and $\Prob(j\mid y)$ are the PIPs of $x_i$ and
%	$x_j$.  $J_{\mathit{LS}}$ can be interpreted as the posterior odds
%	ratio of the models including both \textit{i} and \textit{j} vs the
%	models that include them only individually (see
%	\citealt{LeySteel2007}).
	\item Another jointness measure was proposed in
		\cite{Strachan-JAE-2009} as
		(\lstinline!"Strachan Measure"! in our package):
		\begin{equation}
			J_{\mathit{St}}(A,B) = \Prob(A\mid y)\Prob(B\mid y)\ln\left(
			\frac{\Prob(AB\mid y)}{\Prob(A\overline{B}\mid y)\Prob(\overline{A}B\mid y)}\right)
			\sim J_{\mathit{ST}} = \frac{a+b}{N}\frac{a+c}{N}\ln\left(\frac{aN}{bc}\right).
			\label{St_jointness}
		\end{equation}
	\item Next jointness measure is Yule's Q defined as
		(\lstinline!"Yule's Q Measure (YQ)"! in our package):
		\begin{equation}
			J_{\mathit{YQ}}(A,B) = 
			\frac{\Prob(AB)\Prob(\overline{AB})-\Prob(A\overline{B})\Prob(\overline{A}B)}
			{\Prob(AB)\Prob(\overline{AB})+\Prob(A\overline{B})\Prob(\overline{A}B)}
			\sim J_{\mathit{YQ}} = \frac{ad - bc}{ad + bc}.
			\label{Yq_jointness}
		\end{equation}
		\item The last jointness measure (the default) is
			modified Yule's Q defined as
			(\lstinline!"Modified Yule's Q Measure (YQM)"! in our package):
	\begin{equation}		
		J_{\mathit{YQ}} = \frac
		{\left(a+\frac{1}{2}\right)\left(d+\frac{1}{2}\right) -
		\left(b+\frac{1}{2}\right)\left(c+\frac{1}{2}\right)} 
		{\left(a+\frac{1}{2}\right)\left(d+\frac{1}{2}\right) +
			\left(b+\frac{1}{2}\right)\left(c+\frac{1}{2}\right) - \frac{1}{2}}.
		\label{YqM_jointness}
	\end{equation}
\end{enumerate}

According to \cite{DoppelhoferWeeks2009}, we use the classification
of jointness among the variables given in Table~\ref{tab1}.

\begin{table}[hbt!]
	\centering
	\caption{Classification of strength of jointness measures.}\label{tab1}	
	\begin{tabular}{lc}\hline
		Evidence & Jointness statistics \\ \hline
		strong substitutes & $J_{\mathit{LS}}, J_{\mathit{DW}} \le -2 $ \\
		significant substitutes & $-2 <J_{\mathit{LS}}, J_{\mathit{DW}} \le-1 $ \\
		not significantly related & $-1< J_{\mathit{LS}}, J_{\mathit{DW}} <1 $ \\
		significant complements & $ 1\ge J_{\mathit{LS}}, J_{\mathit{DW}} <2 $ \\
		strong complements & $ J_{\mathit{LS}}, J_{\mathit{DW}} \ge 2 $ \\ \hline
	\end{tabular}	
\end{table} 
